\documentclass[ctcc_full_documentation.tex]{subfiles}
\begin{document}
\chapter{JSON Loaders Documentation}

\section{Introduction}

The \texttt{json\_loaders.py} module loads CTCC configuration from a single combined JSON file when \texttt{CTCC\_INPUT\_MODE=json}. It provides a singleton \texttt{JSONDataSource} that holds the combined JSON in memory and exposes \texttt{get\_data(key)} for section access. A set of \texttt{load\_*} functions mirror \texttt{yaml\_loaders.py}: each maps a JSON key (e.g.\ \texttt{01\_project\_technical\_details}) to the same dataclasses and structures used in YAML mode, so calculation scripts can use the same API regardless of input mode.

\section{Method Overview}

\subsection{Purpose}

\begin{enumerate}
    \item Provide a singleton \texttt{JSONDataSource} that stores the combined JSON; data can be set via \texttt{set\_json\_data()} or auto-loaded from \texttt{CTCC\_JSON\_DATA\_FILE} when \texttt{get\_data()} is first called.
    \item Expose \texttt{get\_data(key)} to return the dict for a given top-level key (e.g.\ \texttt{03\_financing}).
    \item Provide \texttt{load\_*} functions that call \texttt{get\_data()} with the appropriate key and map the dict into the same dataclasses as \texttt{yaml\_loaders} (e.g.\ \texttt{ProjectTechnicalDetails}, \texttt{FinancingDetails}, \texttt{CongestionCurtailmentParams}).
    \item Use \texttt{calculation\_utils} and \texttt{yaml\_loaders} for normalization (e.g.\ \texttt{normalize\_capacity\_mw}, \texttt{normalize\_construction\_type}) and type definitions; \texttt{financial\_utils} for WACC when loading financing.
\end{enumerate}

\subsection{Calculation Approach}

I/O only; no cost/benefit formulas. Keys in the combined JSON match YAML file names (e.g.\ \texttt{01\_project\_technical\_details}, \texttt{17\_congestion\_curtailment\_reductions}). There is no scenario ID in the loader; the combined JSON is assumed to represent one scenario. Scenario selection, if needed, is handled by the caller (e.g.\ server or CLI) before setting the JSON data.

\subsection{Inputs and Outputs}

\textbf{Inputs:}
\begin{itemize}
    \item Combined JSON: set programmatically via \texttt{set\_json\_data(combined\_json)}, or loaded from the file path in \texttt{CTCC\_JSON\_DATA\_FILE} when \texttt{get\_data()} is first called.
    \item Top-level keys in the JSON correspond to YAML section names (e.g.\ \texttt{01\_project\_technical\_details}, \texttt{02\_project\_physical\_details}, \texttt{03\_financing}, \texttt{17\_congestion\_curtailment\_reductions}, \texttt{22\_primary\_bcr\_config}).
\end{itemize}

\textbf{Outputs:}
\begin{itemize}
    \item In-memory dicts or dataclass instances returned to callers (calculation scripts); no files written by this module.
\end{itemize}

\subsection{Integration with Other Scripts}

\begin{itemize}
    \item \texttt{smart\_loaders.py} --- When \texttt{CTCC\_INPUT\_MODE=json}, imports all \texttt{load\_*} and \texttt{get\_data} from \texttt{json\_loaders}; calculation scripts then call \texttt{load\_project\_technical\_details()}, \texttt{load\_financing\_details()}, etc., without knowing the backend.
    \item \texttt{yaml\_loaders.py} --- Supplies dataclass and type definitions (\texttt{ProjectTechnicalDetails}, \texttt{CongestionCurtailmentParams}, \texttt{PhysicalDetailsDetailed}, \texttt{CircuitAndResistanceDetails}, \texttt{FinancingDetails}) and \texttt{normalize\_construction\_type}.
    \item \texttt{calculation\_utils.py} --- \texttt{normalize\_capacity\_mw}, \texttt{get\_converter\_type}.
    \item \texttt{financial\_utils.py} --- \texttt{calculate\_real\_wacc} used in \texttt{load\_financing\_details}.
\end{itemize}

\section{Key Functions}

\subsection{\texttt{JSONDataSource}}

Singleton class that holds the combined JSON. \texttt{get\_data(key)} returns \texttt{self.\_json\_data[key]}; if \texttt{\_json\_data} is \texttt{None}, it first tries \texttt{\_load\_from\_env\_file()} (read from \texttt{CTCC\_JSON\_DATA\_FILE}).

\textbf{Key logic (get\_data):}
\begin{lstlisting}
if self._json_data is None:
    self._load_from_env_file()
if self._json_data is None:
    raise RuntimeError("JSON data not set...")
if key not in self._json_data:
    raise KeyError(f"Key '{key}' not found in JSON data.")
return self._json_data[key]
\end{lstlisting}

\textbf{Key logic (\_load\_from\_env\_file):}
\begin{lstlisting}
json_file_path = os.environ.get("CTCC_JSON_DATA_FILE")
if json_file_path and os.path.exists(json_file_path):
    with open(json_file_path, "r") as f:
        self._json_data = json.load(f)
\end{lstlisting}

\subsection{\texttt{set\_json\_data(combined\_json)}}

Sets the combined JSON dict on the global \texttt{JSONDataSource} instance so all \texttt{load\_*} functions use it.

\subsection{\texttt{load\_project\_technical\_details()}}

Loads \texttt{01\_project\_technical\_details}, normalizes construction type and capacity, and returns \texttt{ProjectTechnicalDetails} (same as \texttt{yaml\_loaders}).

\textbf{Key logic:}
\begin{lstlisting}
project_details = _data_source.get_data("01_project_technical_details")
pd = project_details["project"]
tl = project_details["timeline"]
construction_type = normalize_construction_type(pd["construction_type"])
capacity_mw = normalize_capacity_mw(pd["capacity_mw"])
# ... converter_type, row_agreement_type, etc.
return ProjectTechnicalDetails(
    construction_type=construction_type,
    ac_dc=pd["ac_dc"],
    capacity_mw=capacity_mw,
    ...
)
\end{lstlisting}

\subsection{\texttt{load\_financing\_details()}}

Loads \texttt{03\_financing}, computes real WACC via \texttt{calculate\_real\_wacc}, returns \texttt{FinancingDetails}.

\subsection{\texttt{load\_congestion\_curtailment\_reductions()}}

Loads \texttt{17\_congestion\_curtailment\_reductions}; checks \texttt{01\_project\_technical\_details} for \texttt{reconductoring} to choose \texttt{greenfield\_*} vs \texttt{reconductoring\_*} section. Returns \texttt{CongestionCurtailmentParams} with flow factor, binding hours, prices, curtailment parameters, etc.

\subsection{\texttt{load\_primary\_bcr\_config()}}

Loads \texttt{22\_primary\_bcr\_config} and returns a dict of module enable/disable flags (\texttt{operational}, \texttt{risk}, \texttt{energy}, \texttt{benefits}). If the section is missing or invalid, returns default (all enabled) and optionally warns.

\subsection{Other \texttt{load\_*} functions}

\begin{itemize}
    \item \texttt{load\_physical\_details()} --- \texttt{02\_project\_physical\_details}; returns total miles (sum of \texttt{terrain.terrain\_miles} values).
    \item \texttt{load\_physical\_details\_detailed()} --- returns \texttt{PhysicalDetailsDetailed} with per-terrain miles.
    \item \texttt{load\_circuit\_and\_resistance\_details(category)} --- \texttt{21\_project\_category\_circuit\_and\_resistance\_detail}; returns \texttt{CircuitAndResistanceDetails} for the given category.
    \item \texttt{load\_row\_widths(category)}, \texttt{load\_row\_details()}, \texttt{load\_delay\_costs()}, \texttt{load\_emissions\_details()}, \texttt{load\_contingencies()}, \texttt{load\_environmental\_mitigation()}, \texttt{load\_cost\_timing\_patterns()}, \texttt{load\_afudc\_config()}, \texttt{load\_insurance\_details()}, \texttt{load\_wildfire\_costs()}, \texttt{load\_outage\_costs()}, \texttt{load\_terrain\_data()} --- each calls \texttt{get\_data} with the corresponding key and returns the section dict or derived values.
\end{itemize}

\section{Example}

Combined JSON top-level keys (same as YAML file names):
\begin{lstlisting}
{
  "01_project_technical_details": { "project": {...}, "timeline": {...} },
  "02_project_physical_details": { "terrain": {...} },
  "03_financing": { "financial": {...} },
  "17_congestion_curtailment_reductions": { "greenfield_...": {...}, "reconductoring_...": {...} }
}
\end{lstlisting}

Using the loader (e.g.\ in a calculation script):
\begin{lstlisting}
# Set by caller or auto-load from CTCC_JSON_DATA_FILE
from smart_loaders import load_project_technical_details, load_financing_details
details = load_project_technical_details()
financing = load_financing_details()
\end{lstlisting}

\end{document}
